\documentclass[12pt,a4paper]{article}

\usepackage[utf8]{inputenc}
\usepackage[T1]{fontenc}
\usepackage[french]{babel}
\usepackage{lmodern}
\usepackage{amsmath,amssymb}
\usepackage{geometry}
\usepackage{microtype}
\usepackage{hyperref}
\usepackage{enumitem}
\geometry{top=2.5cm,bottom=2.5cm,left=2.5cm,right=2.5cm}
\hypersetup{hidelinks,pdftitle={Documentation Cortex HCANN}}

\title{Documentation du Bloc Cortex}
\author{Projet HCANN-Agent }
\date{\today}

\begin{document}
\maketitle

\section{Bloc Cortex (\texttt{cortex.py})}

\subsection{Objectif du bloc}
Le bloc cortex regroupe trois sous-modules:
\begin{itemize}
    \item \textbf{MultimodalEncoder}: projection des modalités vision/texte dans un espace sémantique commun;
    \item \textbf{WorkingMemory}: mémoire de travail à slots avec lecture/écriture douce;
    \item \textbf{EntorhinalGateway}: passerelle cortex$\leftrightarrow$hippocampe par attention.
\end{itemize}

\subsection{Inspiration neuroscientifique}
Le bloc cortex est inspiré de plusieurs fonctions attribuées au néocortex et au cortex entorhinal dans la littérature:
\begin{itemize}
    \item \textbf{Néocortex associatif}: intégration de modalités hétérogènes vers une représentation plus abstraite et stable;
    \item \textbf{Mémoire de travail}: maintien transitoire d'un petit nombre d'items pertinents pour le traitement courant;
    \item \textbf{Passerelle entorhinale}: interface fonctionnelle entre représentations corticales et dynamique hippocampique.
\end{itemize}
L'implémentation ne cherche pas à reproduire la biophysique fine des circuits, mais à conserver leur rôle algorithmique (fusion, maintien contextuel, routage vers l'hippocampe).

\subsection{1) Encodeur multimodal}

Soient $\mathbf{v}\in\mathbb{R}^{512}$ (feature vision CLIP) et $\mathbf{t}\in\mathbb{R}^{512}$ (feature texte CLIP).  
Le module applique des projections linéaires:
\begin{align}
\mathbf{z}_v &= \mathrm{Normalize}(W_v\mathbf{v} + b_v), \\
\mathbf{z}_t &= \mathrm{Normalize}(W_t\mathbf{t} + b_t),
\end{align}
avec $W_v, W_t \in \mathbb{R}^{d_s \times 512}$, où $d_s=\texttt{semantic\_dim}$.

La fusion multimodale est une moyenne simple:
\begin{equation}
\mathbf{z} =
\begin{cases}
\mathbf{z}_v & \text{si vision seule,}\\
\mathbf{z}_t & \text{si texte seul,}\\
\frac{1}{2}(\mathbf{z}_v+\mathbf{z}_t) & \text{si vision + texte.}
\end{cases}
\end{equation}

\paragraph{Mode mock}
Quand \texttt{use\_mock\_encoder=True}, les features CLIP réelles sont remplacées par des vecteurs aléatoires normalisés, tout en conservant les dimensions batch et la même interface. Cela accélère les tests sans téléchargement de modèles.

\subsection{2) Working Memory à slots}

On note:
\begin{itemize}
    \item $\mathbf{x}\in\mathbb{R}^{B\times d_w}$ l'entrée batch;
    \item $S\in\mathbb{R}^{N_s\times d_w}$ la matrice des slots (\texttt{wm\_slots} slots).
\end{itemize}

Lecture par attention:
\begin{align}
A &= \frac{\mathbf{x}S^\top}{\sqrt{d_w}},\\
W &= \mathrm{softmax}(A),\\
\mathbf{r} &= WS.
\end{align}

Écriture douce (slot gagnant par échantillon):
\begin{equation}
S_{k^\star} \leftarrow 0.95\,S_{k^\star} + 0.05\,\mathbf{x}_i,
\end{equation}
où $k^\star=\arg\max_k W_{ik}$.

\paragraph{Perte de diversité}
Pour éviter la redondance entre slots:
\begin{equation}
\mathcal{L}_{div}
= \left\| \tilde{S}\tilde{S}^\top - I \right\|_F^2,
\end{equation}
avec $\tilde{S}$ la version normalisée de $S$.

\subsection{3) Passerelle entorhinale}

Le module \texttt{EntorhinalGateway} implémente une attention simplifiée:
\begin{align}
\mathbf{q} &= W_q \mathbf{a},\\
[\mathbf{k},\mathbf{v}] &= \mathrm{split}(W_{kv}\mathbf{m}),\\
\alpha &= \mathrm{softmax}\!\left(\frac{\mathbf{q}\mathbf{k}^\top}{\sqrt{d_k}}\right),\\
\mathbf{h} &= W_o(\alpha \mathbf{v}),
\end{align}
où $\mathbf{a}$ est l'état attracteur cortical et $\mathbf{m}$ l'état de mémoire de travail.

\subsection{Correspondance implémentation}
\begin{itemize}
    \item \texttt{MultimodalEncoder.forward}: extraction CLIP + projection + normalisation + fusion;
    \item \texttt{WorkingMemory.forward}: attention sur slots + mise à jour incrémentale;
    \item \texttt{EntorhinalGateway.forward}: cross-attention cortex$\rightarrow$hippocampe.
\end{itemize}

\subsection{Détail mathématique fonction par fonction}

\paragraph{\texttt{MultimodalEncoder.\_\_init\_\_}}
Initialise les projections:
\[
W_v, W_t \in \mathbb{R}^{d_s \times 512}, \quad b_v,b_t\in\mathbb{R}^{d_s}.
\]
Si \texttt{use\_mock\_encoder=False}, CLIP est chargé puis figé.

\paragraph{\texttt{MultimodalEncoder.\_freeze\_clip}}
Pour chaque paramètre $\theta$ du backbone CLIP:
\[
\theta.\texttt{requires\_grad} \leftarrow \texttt{False}.
\]
En pratique, cela bloque la rétropropagation sur CLIP et réserve l'apprentissage aux têtes de projection.

\paragraph{\texttt{MultimodalEncoder.forward}}
Soit un batch de taille $B$.  
Vision: $\mathbf{V}\in\mathbb{R}^{B\times512}$, texte: $\mathbf{T}\in\mathbb{R}^{B\times512}$.
\begin{align}
\mathbf{Z}_v &= \mathrm{norm}(\mathbf{V}W_v^\top + b_v),\\
\mathbf{Z}_t &= \mathrm{norm}(\mathbf{T}W_t^\top + b_t),\\
\mathbf{Z} &= \frac{1}{n_m}\sum_{m=1}^{n_m}\mathbf{Z}_m,
\end{align}
où $n_m\in\{1,2\}$ est le nombre de modalités présentes.

\paragraph{\texttt{MultimodalEncoder.\_mock\_forward}}
Remplace $\mathbf{V}$ et/ou $\mathbf{T}$ par des gaussiennes:
\[
\tilde{\mathbf{V}},\tilde{\mathbf{T}} \sim \mathcal{N}(0,I),
\]
puis applique les mêmes projections.  
Le code force la cohérence des dimensions batch en alignant la taille texte sur la taille image quand les deux sont fournies.

\paragraph{\texttt{WorkingMemory.forward}}
Entrée $\mathbf{X}\in\mathbb{R}^{B\times d_w}$, slots $S\in\mathbb{R}^{N_s\times d_w}$.
\begin{align}
A &= \frac{\mathbf{X}S^\top}{\sqrt{d_w}} \in \mathbb{R}^{B\times N_s},\\
W &= \mathrm{softmax}(A),\\
\mathbf{R} &= WS \in \mathbb{R}^{B\times d_w}.
\end{align}
Écriture (sans gradient):
\[
S_{k^\star(i)} \leftarrow 0.95\,S_{k^\star(i)} + 0.05\,\mathbf{X}_i.
\]

\paragraph{\texttt{WorkingMemory.diversity\_loss}}
Si $\hat{S}$ est la version normalisée de $S$:
\[
\mathcal{L}_{div} = \left\| \hat{S}\hat{S}^\top - I \right\|_F^2.
\]
Le terme pénalise les slots colinéaires et favorise la spécialisation.

\paragraph{\texttt{EntorhinalGateway.forward}}
Avec état attracteur $\mathbf{a}$ et état mémoire de travail $\mathbf{m}$:
\begin{align}
\mathbf{q} &= W_q\mathbf{a},\\
[\mathbf{k},\mathbf{v}] &= \mathrm{split}(W_{kv}\mathbf{m}),\\
\alpha &= \mathrm{softmax}\!\left(\frac{\mathbf{q}\mathbf{k}^\top}{\sqrt{d_k}}\right),\\
\mathbf{h} &= W_o(\alpha\mathbf{v}).
\end{align}
La sortie $\mathbf{h}$ est projetée dans l'espace hippocampique (\texttt{hpc\_size}).

\subsection{Interface détaillée des fonctions}

\paragraph{\texttt{MultimodalEncoder.forward(images=None, texts=None)}}
\begin{itemize}[leftmargin=*]
    \item \textbf{Entrées}
    \begin{itemize}
        \item \texttt{images}: \texttt{torch.Tensor} de forme $(B, C, H, W)$, optionnel;
        \item \texttt{texts}: liste de $B$ chaînes (ou structure équivalente), optionnel.
    \end{itemize}
    \item \textbf{Sortie}
    \begin{itemize}
        \item embedding sémantique \texttt{torch.Tensor} de forme $(B, d_s)$.
    \end{itemize}
    \item \textbf{Erreur métier}
    \begin{itemize}
        \item si aucune modalité n'est fournie, l'implémentation lève une exception.
    \end{itemize}
\end{itemize}

\paragraph{\texttt{WorkingMemory.forward(x)}}
\begin{itemize}[leftmargin=*]
    \item \textbf{Entrée}: \texttt{x} de forme $(B,d_w)$.
    \item \textbf{Sorties}:
    \begin{itemize}
        \item \texttt{read}: $(B,d_w)$, contenu lu depuis les slots;
        \item \texttt{w}: $(B,N_s)$, poids d'attention sur slots.
    \end{itemize}
\end{itemize}

\paragraph{\texttt{WorkingMemory.diversity\_loss()}}
\begin{itemize}[leftmargin=*]
    \item \textbf{Entrée}: aucune (utilise les slots internes).
    \item \textbf{Sortie}: scalaire \texttt{torch.Tensor} ($\mathcal{L}_{div}$).
\end{itemize}

\paragraph{\texttt{EntorhinalGateway.forward(attractor\_state, wm\_state)}}
\begin{itemize}[leftmargin=*]
    \item \textbf{Entrées}
    \begin{itemize}
        \item \texttt{attractor\_state}: état attracteur cortical;
        \item \texttt{wm\_state}: état mémoire de travail.
    \end{itemize}
    \item \textbf{Sortie}: projection vers l'espace hippocampique, dimension finale \texttt{hpc\_size}.
\end{itemize}

\subsection{Traçage calculatoire (formules $\leftrightarrow$ code)}
\begin{enumerate}[leftmargin=*]
    \item \textbf{Extraction des features} : CLIP (vision/texte) ou vecteurs mock aléatoires.
    \item \textbf{Projection linéaire} : multiplication affine par \texttt{vision\_proj} et \texttt{text\_proj}.
    \item \textbf{Normalisation} : \texttt{F.normalize(..., dim=-1)} pour stabiliser l'échelle.
    \item \textbf{Fusion} : moyenne arithmétique des modalités présentes.
    \item \textbf{Mémoire de travail} : calcul attention $A$, poids $W$, lecture $R=WS$.
    \item \textbf{Écriture slot} : mise à jour EMA du slot gagnant (\texttt{argmax} sur $W$).
\end{enumerate}

\end{document}
